\documentclass{article}
\usepackage{../techreport/arxiv}

\usepackage[utf8]{inputenc}
\usepackage[T1]{fontenc}
\usepackage{hyperref}
\usepackage{url}
\usepackage{booktabs}
\usepackage{amsmath}
\usepackage{amsfonts}
\usepackage{array}
\usepackage{microtype}
\usepackage{siunitx}
\usepackage{enumitem}

\newcommand{\CurrentLogRows}{9,636}
\newcommand{\CurrentLogRunCount}{2,417}
\newcommand{\CurrentLogTransitions}{2,408}
\newcommand{\PseudoLabelExamples}{5,513}
\newcommand{\PseudoLabelSwitchCount}{381}
\newcommand{\FutureCandidateCount}{2,370}
\newcommand{\FutureCandidatePct}{98.4}
\newcommand{\FutureTopTwoCount}{1,706}
\newcommand{\FutureTopTwoPct}{70.8}
\newcommand{\FutureRankedBetterCount}{1,269}
\newcommand{\FutureRankedBetterPct}{52.7}
\newcommand{\ChangedPseudoTopOneCount}{209}
\newcommand{\ChangedPseudoTopTwoCount}{272}
\newcommand{\ChangedMiniRecoverCount}{166}
\newcommand{\ChangedBothAgreeMissCount}{195}
\newcommand{\BenchmarkTrainExamples}{3,307}
\newcommand{\BenchmarkValExamples}{1,103}
\newcommand{\BenchmarkTestExamples}{1,103}

\newcommand{\ReplayAllLogs}{10}
\newcommand{\ReplayAllLogFixes}{11,351}
\newcommand{\ReplayPseudoExamples}{6,529}
\newcommand{\ReplayChangedExamples}{453}
\newcommand{\ReplayAccuracy}{95.99}
\newcommand{\ReplayChangedRecall}{72.19}
\newcommand{\ReplayUnchangedAccuracy}{97.76}
\newcommand{\ReplayThreeWayGateActivations}{3,118}

\newcommand{\ModelBenchmarkRows}{%
Current final matcher & 90.66\% & 0.00\% & 100.00\% \\%
Mini-HMM else heuristic & 92.57\% & 34.95\% & 98.50\% \\%
Lowest-distance baseline & 90.75\% & 54.37\% & 94.50\% \\%
Two-way tree gate & 93.56\% & 46.60\% & 98.40\% \\%
Three-way tree gate & 93.11\% & 42.72\% & 98.30\% \\%
Three-way fixed-lag smoother & 94.20\% & 55.34\% & 98.20\% \\%
Tree candidate ranker & 93.11\% & 50.49\% & 97.50\% \\%
Logistic ranker & 91.02\% & 41.75\% & 96.10\% \\%
Tiny MLP ranker & 87.04\% & 20.39\% & 93.90\% \\%
Online linear ranker & 90.75\% & 48.54\% & 95.10\% \\
}


\title{Fixed-Lag Self-Supervised Personalization of Lightweight Map Matching for Mobile Speed-Limit Assistance}

\author{
  Raphael Volz \\
  Pforzheim University \\
  \texttt{raphael.volz@hs-pforzheim.de}
}

\begin{document}
\maketitle

\begin{abstract}
Smartphone-based intelligent speed assistance needs map matching that remains stable under noisy consumer GPS, intermittent signal loss, and strict on-device latency budgets. The deployed YouSpeed runtime therefore uses a lightweight matcher built from four information layers: single-fix geometry, fix-sequence continuity state, single-way topology, and way-sequence priors including a speed-dependent switch-feasibility term, rather than a full route-level probabilistic model. This paper studies whether rolling driving logs can personalize that matcher without cloud supervision or dense ground-truth labeling. We first decompose the problem by information items available from single fixes, fix sequences, single ways, and way sequences, and then evaluate a fixed model ladder derived from those layers. We propose a fixed-lag self-supervision scheme in which later stable fixes provide hindsight pseudo-labels for earlier ambiguous candidate sets. On a retained multi-log benchmark with \CurrentLogRows{} matched fixes and \CurrentLogTransitions{} selected-way transitions after dropping strict subset logs, the future-selected road was already present in the previous candidate set in \FutureCandidateCount{} transitions (\FutureCandidatePct{}\%), and already in the previous top-2 in \FutureTopTwoCount{} transitions (\FutureTopTwoPct{}\%). Under a strict five-fix hindsight rule, the benchmark yields \PseudoLabelExamples{} pseudo-labeled examples, including \PseudoLabelSwitchCount{} true switch cases. On the test split, the current final matcher is maximally conservative at 90.66\% accuracy with 100.00\% keep-current accuracy but 0.00\% switch-case recall. The best offline result now comes from a three-way fixed-lag smoother over the expert gate at 94.20\% accuracy with 55.34\% switch-case recall and 98.20\% keep-current accuracy, while the lowest-distance baseline already recovers 54.37\% of switch cases. A separate end-to-end replay across all \ReplayAllLogs{} available logs covers \ReplayPseudoExamples{} hindsight-labeled examples and reaches \ReplayAccuracy{}\% overall accuracy, \ReplayChangedRecall{}\% switch-case recall, and \ReplayUnchangedAccuracy{}\% keep-current accuracy, exposing the closed-loop gap between offline correction potential and runtime self-generated context. The contribution is a deployable personalization formulation, an information-item-based experimental design, and a fixed evaluation protocol for mobile ITS.
\end{abstract}

\keywords{intelligent transportation systems \and intelligent speed assistance \and map matching \and self-supervision \and on-device learning \and personalization}

\section{Introduction}
Intelligent speed assistance (ISA) is mandatory for new vehicle types in the European Union, but the installed vehicle base is far larger than annual new registrations \cite{eu2019_2144,eu2021_1958,eu_mandatory_safety_2024}. This creates a practical retrofit opportunity for smartphone-based assistance that can infer speed context from open map data under intermittent connectivity and consumer-grade sensing.

The deployed YouSpeed app is designed for this product setting. Its speed-limit inference pipeline runs entirely on-device against a local OpenStreetMap-derived database and must therefore keep matching logic computationally small and auditable. In contrast to full Hidden Markov Model (HMM) or path-inference pipelines, the current runtime uses per-fix geometry, continuity-aware hysteresis, same-reference segment promotion, a speed-dependent switch-feasibility prior, tunnel gating, and a small auxiliary sequence check. This design is product-rational, but it creates an obvious question: can repeated driving on the same routes personalize heuristic weights without giving up the low-latency architecture?

This paper studies that question for an ITS audience focused on deployable mobile driver assistance rather than offline leaderboard accuracy. The central hypothesis is that future stable fixes can provide delayed supervision for earlier ambiguous matches. The study design is fixed in advance and grounded in first principles: compare models by which information items they use from single fixes, fix sequences, single ways, and way sequences, while the legal speed-limit derivation remains deterministic and auditable.

The paper is organized as a fixed experimental program rather than a sequence of ad hoc model additions. The contributions are:
\begin{enumerate}[leftmargin=*,itemsep=2pt]
\item A precise formulation of fixed-lag hindsight pseudo-labeling for lightweight mobile map matching.
\item A first-principles decomposition of the information items available from iPhone fixes and OSM ways for lightweight matching.
\item A multi-log empirical assessment derived from that decomposition, covering deterministic baselines, shallow expert gates, candidate rerankers, and a fixed-lag sequence smoother under one common protocol.
\item A fixed refresh protocol for adding future logs without changing the analytical setup.
\end{enumerate}

\section{Related Work}
Map matching has long been studied through geometric, topological, and probabilistic sequence formulations \cite{quddus2007}. HMM-based methods and related path-inference approaches remain the standard answer to noisy or sparse trajectories because they use later observations to stabilize earlier candidate choices \cite{newson2009,lou2009,goh2012,he2013,hunter2014,quddus2015,jagadeesh2017,che2018}. This is the conceptual backdrop for our work: later fixes contain useful information.

Recent work also shows that learning-based map matching is viable when labels are limited. Jin et al.\ train a Transformer map-matching model with transfer learning to reduce dependence on expensive ground-truth data \cite{jin2022transformer}. Safarzadeh and Wang use deep inverse reinforcement learning with multi-intention modeling to handle low-sampling trajectories and driver heterogeneity \cite{safarzadeh2024midirl}. Meng and Zhao push further toward adaptation under limited data with a meta-learning framework that combines user-specific and region-shared patterns for online map matching \cite{meng2026mpm}.

Personalization from historical trajectories is also established outside narrow map-matching formulations. Sun et al.\ show that a shared-plus-individual linear adaptation model improves personalized route-choice prediction as more user data becomes available \cite{sun2023humancentric}. Zhang et al.\ integrate meta-learning into inverse-reinforcement route planning for sparse trajectory coverage \cite{zhang2024metairl}. Most closely related, Zhang et al.\ incorporate personalized road-selection preferences directly into an HMM matcher \cite{zhang2025pphmm}.

Our contribution is narrower than these lines of work. We do not propose a new full-sequence model. Instead, we ask how a lightweight, already deployed smartphone matcher can use fixed-lag hindsight from its own rolling logs to personalize candidate ranking on device. The novelty claim is therefore in formulation and systems fit: self-supervised adaptation for a low-latency mobile ISA stack.

\section{Information Items for Lightweight Mobile Matching}
Lightweight mobile map matching can be decomposed by what information is available from a single iPhone fix, what only appears across a sequence of fixes, what is available from a single way, and what only appears across a sequence of ways. That decomposition is useful because it turns the model-design question into an information-allocation question: which information items does each model family consume, and which additional information layer does a more complex model actually buy?

\subsection{Single-fix information from iPhone sensing}
Core Location delivers location updates through \texttt{CLLocationManager} and exposes location observations as \texttt{CLLocation} objects \cite{apple_cllocationmanager,apple_cllocation}. At the API level, a single fix can contain coordinate, altitude, horizontal and vertical accuracy, speed, speed accuracy, course, course accuracy, timestamp, and source metadata. iOS can additionally provide heading observations as \texttt{CLHeading}, including true heading, magnetic heading, and heading accuracy \cite{apple_clheading}. In other words, a smartphone fix is not just a point; it is a point plus motion and confidence metadata.

The current iPhone client uses the following single-fix items in the matcher: coordinate, timestamp, speed, course, course accuracy, horizontal accuracy, and vertical accuracy. The location manager is configured for automotive navigation with best available accuracy and a \SI{10}{m} distance filter. The current runtime does \emph{not} consume separate heading updates, speed accuracy, or source metadata, even though those items exist in the Apple API. The gate experiments additionally use a small application-level GPS quality bin derived from these raw accuracy signals rather than a separate sensor channel.

\subsection{Information derived from fix sequences}
Once multiple fixes are available, new information appears that does not exist in a single observation: elapsed time, displacement, speed and course evolution, persistence of candidate support, abrupt GPS-quality changes, and stable future runs. In the deployed app, the fix-sequence layer is represented through rolling continuity state such as recent GPS signal loss, tunnel mode, rolling matched-way history, and a small auxiliary hypothesis list. In the learning pipeline, the same layer yields fixed-lag hindsight pseudo-labels and the candidate-stability windows used for evaluation.

\subsection{Information from single ways}
On the map side, a single way contributes its own local evidence. The current OSM-derived database exposes polyline geometry, bounding box, approximate heading, endpoint proximity, street name, street reference, highway class, service and tunnel tags, and speed-limit tags. These are single-way items because they can be evaluated without knowing what happened on the previous fix or previous road.

\subsection{Information from way sequences}
Way-sequence information only appears when ways are related over time or topology. In the current matcher this includes same-way continuity, same-reference corridor continuity, linked-way continuity, recent matched-way history, recent street-reference history, rolling auxiliary hypotheses, and tunnel-chain continuity. These are the information items behind preferred-way hysteresis, same-reference promotion, linked-way fallback, recent-way fallback, and tunnel gating.

\begin{table}[!htbp]
\centering
\caption{Basic information items available to lightweight mobile map matching in the current iPhone and map stack.}
\label{tab:information_items}
\footnotesize
\setlength{\tabcolsep}{4pt}
\renewcommand{\arraystretch}{1.06}
\begin{tabular}{>{\raggedright\arraybackslash}p{2.1cm}>{\raggedright\arraybackslash}p{5.0cm}>{\raggedright\arraybackslash}p{5.9cm}}
\toprule
Layer & Information items & Matching use \\
\midrule
Single fix & Coordinate, timestamp, horizontal/vertical accuracy, speed, speed accuracy, course, course accuracy, optional source metadata and heading & Candidate retrieval radius, point-to-way distance, heading penalty, speed-gated turn feasibility, logging, and confidence gating \\
Fix sequence & Delta time, displacement, speed/course evolution, GPS-loss episodes, persistence of candidate support, stable future runs & Continuity state, tunnel mode, recent-loss protection, auxiliary hypotheses, hindsight pseudo-labels, and fixed-lag smoothing experiments \\
Single way & Polyline geometry, bounding box, approximate heading, endpoint proximity, street name, ref, highway/service/tunnel tags, speed-limit tags & Geometry score, endpoint specialist, heading comparison, speed-limit derivation, and tunnel candidate filtering \\
Way sequence & Same way ID, same ref corridor, linked ways, recent way history, recent ref history, tunnel chains, rolling hypotheses & Preferred-way hysteresis, same-ref promotion, linked-way continuity, mini-HMM transitions, and tunnel continuity \\
\bottomrule
\end{tabular}
\end{table}

This decomposition is the organizing principle of the paper. The runtime is first described in terms of how it uses these information layers, and the benchmark is then built as a planned ladder that either removes or adds one layer at a time.

\section{From Information Items to the Experimental Program}
\subsection{Current runtime matcher}
The implemented matcher is deliberately simpler than full route-level inference. For each GPS fix \(t\), it retrieves candidate ways \(i \in C_t\) and computes a geometry-driven score
\[
s_{\mathrm{geom}}(t,i)=
\begin{cases}
d_{t,i} + \lambda_h \cdot \Delta\theta_{t,i}, & \text{if course is reliable} \\
d_{t,i}, & \text{otherwise}
\end{cases}
\]
where \(d_{t,i}\) is point-to-way distance and \(\Delta\theta_{t,i}\) is bidirectional heading mismatch between the observed course and the way's approximate heading. The final selection then injects way-sequence information: preferred-way hysteresis, same-reference promotion, linked-way continuity, recent-way fallback, tunnel gating, and a small auxiliary sequence check. In addition, the runtime applies a speed-dependent switch-feasibility prior: sharp high-speed transitions are penalized when the current way remains well aligned with the observed heading, while low-speed connected transitions are supported only when heading evidence improves toward the candidate. The runtime also records a continuity-aware trace score that lexicographically prioritizes continuity class over raw geometry within each fix.

This structure matters for learning because it already exposes a clean separation of information layers. Geometry and tags supply single-fix and single-way evidence. Preferred-way, same-ref, linked-way, recent-way, and tunnel state supply way-sequence evidence. The auxiliary sequence component contributes a compact hypothesis history. The learned problem is therefore usually not ``discover any road from scratch.'' It is ``reorder already available candidates given a richer use of the same information items.'' 

\subsection{Fixed-lag hindsight labeling}
Let the matcher output a selected way \(\hat{y}_t\) and an ordered candidate set \(C_t\) at fix \(t\). For a future window of length \(K\), define the hindsight pseudo-label \(y_t^\star\) as the majority way among the next \(K\) selected outputs, subject to three constraints:
\begin{enumerate}[leftmargin=*,itemsep=2pt]
\item the majority agreement must exceed a threshold,
\item the majority way must already be present in \(C_t\),
\item the majority way must persist for a minimum consecutive run length.
\end{enumerate}
This turns future trajectory consistency into delayed supervision without requiring external annotation. In this paper we use a strict setting of \(K=5\), minimum agreement ratio \(0.8\), and minimum future run length \(5\). The critical point is that the future window contributes \emph{fix-sequence} information; it does not invent new way information.

\subsection{Experiments derived from information layers}
The benchmark is planned as an information ladder with four questions fixed before model comparison:
\begin{enumerate}[leftmargin=*,itemsep=2pt]
\item How much matching quality is already available from single-fix and single-way information alone?
\item How much additional stability comes from way-sequence continuity and compact recent-history state?
\item Given the same candidate set and continuity state, is explicit expert arbitration or candidate reranking a better use of the available information?
\item Does future fix-sequence information at inference time add validation-safe value beyond pointwise expert arbitration?
\end{enumerate}

The deterministic baselines answer the first two questions. The lowest-distance baseline uses only single-fix position and single-way geometry. The current final matcher adds way-sequence continuity and the hand-coded switch-feasibility prior. The mini-HMM arbitration baseline adds compact hypothesis transitions on top of that.

The expert-gate models answer the third question while holding the state space small and auditable. The two-way gate arbitrates between the current selected way and the lowest-distance candidate. The three-way gate adds an endpoint specialist. Both models stay in the single-fix plus way-sequence regime: they do not consume explicit future fixes at inference time.

The candidate-ranker models ask the same third question from a different angle: given the same candidate set, can nonlinear scoring exploit the available single-fix and way-sequence evidence better than hand-tuned heuristics? Following the pointwise view of learning to rank \cite{liu2009ltr}, each candidate \(i \in C_t\) is represented by a feature vector \(x_{t,i}\) and scored by a function \(f(x_{t,i})\). At inference time, the model selects
\[
\tilde{y}_t = \arg\max_{i \in C_t} f(x_{t,i}),
\]
where the hindsight pseudo-label candidate is the positive training target and all other candidates in the same set are negative. The candidate-ranker family therefore adds model capacity, not new sensor channels.

The fixed-lag smoother answers the fourth question. It is the only benchmark model that consumes additional \emph{future fix-sequence information at inference time}. It retains the three-way gate as its emission model and applies a short online-safe dynamic program over a horizon \(L\):
\[
\tilde{y}_{t:t+L}
=
\arg\max_{z_{t:t+L}}
\sum_{\tau=t}^{t+L} \log p_{\mathrm{gate}}(z_\tau \mid x_\tau)
\;+\;
\sum_{\tau=t+1}^{t+L} \psi(z_{\tau-1}, z_\tau),
\]
where \(p_{\mathrm{gate}}\) is the three-way gate probability over expert states and \(\psi\) is a small transition model with bonuses for staying on the same way or reference and penalties for implausible switches. Only the oldest decision in the fixed-lag window is emitted, so the formulation remains compatible with online operation.

\begin{table}[!htbp]
\centering
\caption{Information families used by the planned model ladder.}
\label{tab:model_information_matrix}
\scriptsize
\setlength{\tabcolsep}{3.5pt}
\renewcommand{\arraystretch}{1.05}
\begin{tabular}{>{\raggedright\arraybackslash}p{2.9cm}>{\raggedright\arraybackslash}p{1.8cm}>{\raggedright\arraybackslash}p{1.7cm}>{\raggedright\arraybackslash}p{2.35cm}>{\raggedright\arraybackslash}p{2.45cm}}
\toprule
Model & Single fix & Fix sequence & Single way & Way sequence \\
\midrule
Lowest-distance baseline & Position & -- & Geometry distance & -- \\
Current final matcher & Position, course, speed, quality & Recent loss flag & Geometry, axis heading, endpoints, tags & Continuity state and turn prior \\
Mini-HMM arbitration & Position, course, speed, quality & Rolling hypotheses & Geometry, endpoints, tags & Hypothesis transitions over recent way/ref state \\
Two-way tree gate & Speed, horiz. acc., quality bin & -- & Current versus nearest distance, endpoint, class & Continuity-derived current expert \\
Three-way tree gate & Speed, horiz. acc., quality bin & -- & Expert outputs over current, nearest, endpoint candidates & Continuity-derived expert state \\
Tree candidate ranker & Speed, course, quality & -- & Candidate distance, endpoint, tags & Continuity and expert-relative deltas \\
Logistic / tiny MLP / online linear rankers & Speed, course, quality & -- & Candidate distance, endpoint, tags & Continuity class, heuristic match, mini-HMM match \\
Three-way fixed-lag smoother & Gate emissions & Short future window & Expert-state attributes & Same-way/ref/highway transitions \\
\bottomrule
\end{tabular}
\end{table}

Table~\ref{tab:model_information_matrix} should be read as the planned experimental program: each model family answers a different question about how much value comes from adding or reallocating one information layer while keeping the remaining layers fixed.

Within the candidate-ranker family, the linear and neural baselines are standard and are included for calibration rather than novelty. Logistic regression maps a linear score to a Bernoulli probability through the sigmoid link \cite{cox1958binary},
\[
f_{\mathrm{log}}(x)=\sigma(w^\top x + b)=\frac{1}{1+\exp(-(w^\top x+b))},
\]
while the tiny multilayer perceptron augments the linear score with one hidden layer and nonlinear activation \cite{rumelhart1986backprop},
\[
f_{\mathrm{mlp}}(x)=v^\top \phi(Wx+b_1)+b_2.
\]
The online linear ranker uses the same candidate information items but updates the weight vector sequentially after each hindsight-labeled example. It should not be read as a new algorithmic family. In this paper it is a lightweight engineering prototype inspired by classic online linear learners such as the perceptron \cite{rosenblatt1958perceptron} and passive-aggressive methods \cite{crammer2006pa}. The novelty claim of the paper lies not in these standard model classes themselves, but in how fixed-lag hindsight pseudo-labels are used to personalize a deployed mobile matcher under strict runtime and privacy constraints.

\section{Multi-Log Study}
\subsection{Data and pseudo-label yield}
The retained offline benchmark uses seven non-subset drive-match logs with \CurrentLogRows{} matched fixes, \CurrentLogRunCount{} runs of selected ways, and \CurrentLogTransitions{} transitions. Three additional logs in the repository are strict subsets or duplicates of these retained logs and are excluded from the train/validation/test benchmark to avoid duplication, but are included in the separate end-to-end replay validation. The retained logs are transformed into pseudo-label examples by fixed-lag hindsight labeling and then exported into gate and candidate-ranking datasets.

The hindsight signal is strong. In \FutureCandidateCount{} of \CurrentLogTransitions{} transitions (\FutureCandidatePct{}\%), the road chosen next was already present in the previous fix's candidate set. In \FutureTopTwoCount{} transitions (\FutureTopTwoPct{}\%), it was already in the previous top-2. Moreover, the future road was already ranked better than the current selected road in \FutureRankedBetterCount{} transitions (\FutureRankedBetterPct{}\%). Under the strict five-fix labeling rule, the log yields \PseudoLabelExamples{} pseudo-labeled examples, of which \PseudoLabelSwitchCount{} are true switch cases where hindsight prefers a different road than the current final output.

Among those \PseudoLabelSwitchCount{} switch cases, the hindsight label is already rank-1 in \ChangedPseudoTopOneCount{} cases and top-2 in \ChangedPseudoTopTwoCount{} cases. This confirms that the task is usually not missing-road discovery. The dominant problem class is earlier selection of a plausible alternative already present in the candidate set.

\subsection{Model comparison}
The benchmark uses a chronological train/validation/test split with \BenchmarkTrainExamples{}/\BenchmarkValExamples{}/\BenchmarkTestExamples{} examples. Reported test metrics are overall accuracy, recall on switch cases, and \emph{keep-current accuracy}: accuracy restricted to cases where hindsight agrees with the current selected way and the model should therefore avoid an unnecessary switch. The main table is organized as an information-layer comparison: each row corresponds to a different allocation of single-fix, fix-sequence, single-way, and way-sequence evidence.

\begin{table}[!htbp]
\centering
\caption{Chronological test-set results on hindsight pseudo-labels from the retained multi-log benchmark.}
\label{tab:model_results}
\begin{tabular}{lrrr}
\toprule
Model & Accuracy & Switch recall & Keep-current accuracy \\
\midrule
\ModelBenchmarkRows
\bottomrule
\end{tabular}
\end{table}

Four observations matter most for the information-layer interpretation.

First, the current final matcher is conservative: it preserves keep-current cases perfectly in this benchmark, but recovers none of the switch cases. Second, the lowest-distance baseline is already very strong, recovering 54.37\% of switch cases at 90.75\% overall accuracy. Any learned method must therefore beat a nontrivial geometric baseline, not only the current heuristic. Third, on the refreshed retained corpus the best offline result now comes from the three-way fixed-lag smoother, which reaches 94.20\% accuracy, 55.34\% switch-case recall, and 98.20\% keep-current accuracy. The gain over the pointwise gate is still incremental, because the smoother changes only 17 of 1{,}103 test windows, but it is no longer a validation no-op. Fourth, the closed-loop replay across all \ReplayAllLogs{} logs remains the harder deployment-relevant benchmark: offline correctable ambiguity does not translate one-for-one into recurrent runtime behavior.

\subsection{Error structure and model allocation}
The switch cases split cleanly by model family. In \ChangedMiniRecoverCount{} cases, the auxiliary sequence logic already matches the hindsight pseudo-label when it disagrees with the heuristic. These are arbitration cases inside the existing expert set. In another \ChangedBothAgreeMissCount{} cases, both heuristic and auxiliary sequence logic agree on the wrong road. These are reranking cases, because a simple expert gate cannot recover them. The high values of \ChangedPseudoTopOneCount{} rank-1 and \ChangedPseudoTopTwoCount{} top-2 hindsight labels confirm that most corrections are reordering problems inside the existing candidate set, not missing-road discovery. The fixed-lag smoother improves over the pointwise gate by adjusting only a small minority of windows, so most recoverable ambiguity is still allocated to expert arbitration and candidate reranking rather than large sequence-level corrections.

\subsection{Closed-loop replay across all available logs}
The offline hindsight benchmark still uses logged candidate sets and logged trajectory context. To quantify the gap to deployment, we added a closed-loop replay test that replays all \ReplayAllLogs{} available logs in order, rebuilds matcher state from the runtime's own outputs, and then evaluates against the same five-fix hindsight labels. This replay covers \ReplayAllLogFixes{} fixes and \ReplayPseudoExamples{} hindsight-labeled examples. It reaches \ReplayAccuracy{}\% overall accuracy, \ReplayChangedRecall{}\% switch-case recall, and \ReplayUnchangedAccuracy{}\% keep-current accuracy, while activating the deployed three-way gate \ReplayThreeWayGateActivations{} times.

This result is important because it is much harsher than the offline benchmark. Early deviations alter later context, so local correction potential does not translate directly into closed-loop switch recovery. Scientifically, this means the hindsight benchmark should be read as an estimate of \emph{correctable ambiguity}, whereas the replay test measures the stability of the full recurrent system. Both views are needed.

\section{Discussion and Evaluation Protocol}
The multi-log offline benchmark and the all-log replay together support the scientific validity of the idea, but they do not yet prove broad generalization. Future fixes do not provide absolute truth; they provide high-confidence pseudo-labels. This remains vulnerable to repeated ambiguity on parallel roads, tunnels, missing-map segments, or long periods of bad GPS. The method is therefore only defensible when training is restricted to high-confidence hindsight windows.

The strongest research claim under this design is not ``deep learning for map matching.'' It is narrower and more credible: lightweight self-supervised personalization can tune a deployed heuristic matcher using fixed-lag hindsight from rolling logs, and the fixed model comparison places the present gains in strong geometric baselines, shallow expert arbitration, and a modest amount of fixed-lag smoothing on the retained offline split. For an ITS audience, this is attractive because it preserves operational constraints: low-latency inference, offline capability, interpretable fallbacks, and privacy-preserving local adaptation.

The evaluation protocol is fixed in advance:
\begin{enumerate}[leftmargin=*,itemsep=2pt]
\item ingest a new drive-match log and regenerate hindsight pseudo-labels,
\item refresh gate and candidate datasets under the same strict labeling rule,
\item rerun the baseline, expert-gate, candidate-ranker, and fixed-lag smoother benchmark,
\item rerun the closed-loop replay test across all available logs,
\item update only the generated LaTeX fragments that feed this manuscript,
\item compare warm-up behavior, switch-case recall, and keep-current stability across logs, drivers, and route strata.
\end{enumerate}

This protocol keeps the paper structure stable while evidence accumulates. It is also specific about what the present data do \emph{not} show: the retained corpus now supports a small but real fixed-lag smoothing gain, yet the closed-loop replay still shows that offline improvements do not directly certify deployment behavior. Additional route types such as tunnels and high-speed motorway corridors can therefore be added under the same protocol without changing the analytical framing.

\section{Conclusion}
This paper documented a deployed lightweight mobile matcher for offline speed-limit assistance and evaluated a fixed-lag self-supervision strategy for personalizing it from rolling driving logs. The retained multi-log benchmark shows that future stable fixes frequently reveal which earlier candidate should have been selected. Under the refreshed model comparison, the current final matcher remains highly conservative, the lowest-distance baseline is already hard to beat, and the best offline result now comes from a three-way fixed-lag smoother over the expert gate. The closed-loop replay test shows that offline hindsight gains do not transfer one-for-one into recurrent runtime behavior. The main result is therefore not architectural reinvention but a structured evaluation framework in which additional logs, users, and route types can be assessed under the same offline and closed-loop criteria.

\bibliographystyle{unsrtnat}
\bibliography{../share/IEEEabrv,../share/references}

\end{document}
